\cotefile{29}
\scanpage{24}
% archive sheet number: 23
telle que pour tout objet $M$ d'un $\underline{C}_{S'}$, il existe un recouvrement de $S'$ par des $S'_i$ tel que pour tout $i$, $M_{S'_i}$ soit \deleted{dominé} majo\unsure{r}é par un $M_i$ objet de $\underline{C}'_{S'_i}$, \deleted{- alors} et supposant que le foncteur réalisation géométrique pour $\underline{R}'$ est celui induit pour le foncteur réal. géom. pour $\underline{R}$, - alors on constate que $\underline{R}$ et $\underline{R}'$ sont ``équivalents''. Comme cas particulier, signalons l'exemple de la donnée kummérienne a), dans le cas où les $D_i$ admettent des équations globales par des sections $c_i$ de $\underline{O}_S$. Alors on peut considérer le systême projectif des $Z_{\underline{c}}^{\underline{n}}$, pour $\underline{n}=(n_i)_{i\quad I}$ variables, premiers \deleted{aux \unsure{caractéristiques|caractéritiques} résiduelles} à la caractéristique de $S$ (supposée exister). Considérons de plus le systême projectif des groupes $\underline{Z}_{\underline{n}}=\quad\underline{Z}/n_i\underline{Z}$, et supposons donné un isomorphisme entre celui-ci et le systême projectif des ${}_{\underline{n}}$ (il en existe lorsque $S$ est un schéma sur un corps algébriquement clos, par exemple). On peut alors considérer les $Z_{\underline{c}}^{\underline{n}}$ comme des schémas à groupes $\underline{Z}_{\underline{n}}$, formant un systême projectif, \deleted{d'où} donc on est sous les conditions de l'exemple 2, définissant une donnée de ramification $\underline{R}'$ sur $S$. On constat alors que cette donnée est équivalente à la donnée $\underline{R}$ décrite dans le cas particulier kummérien a). Comme autre exemple de données de ramification équivalentes\deleted{\unsure{:|;|.}}, signalons le suivant, où $\underline{R}$ est la donnée de ramification considérée dans l'exemple 1. On se donne un revêtement principal $P$ de $S$, de groupe $\unsure{H|\deleted{H}G|\deleted{G}H}'$ (opérant à droite), et on suppose que $G$ opère dessus librement à gauche, en commutant aux opérations de $G'$. Soit alors $Z'=Zx^{G}P$, considéré comme revêtement à groupe d'opérateurs $G'$, d'où une donnée de ramification $\underline{R}'$. Cette dernière est équivalente à $\underline{R}$, comme on voit en considérant l\deleted{es}'équivalences \deleted{quasi-inverses l'une} de \deleted{l'a\unsure{u}tre} \deleted{entre} $\underline{\mathrm{Rev}}^{\underline{R}}$ \deleted{et}\inserted{dans} ${\unsure{\underline{\mathrm{Rev}}|\mathrm{Rev}}}^{\unsure{R'|\underline{R}'}}$ opérateurs \illegible{3}
